% Slide bảo vệ khóa luận tốt nghiệp.
% Đề tài: Thiết kế kiến trúc tổng thể hạ tầng mạng không dây tại Trường Đại học Đồng Tháp
% Sinh viên: Nguyễn Thị Kỳ Duyên (MSSV: 0023410647 - Lớp: ĐHCNTT23B-CS)
% GVHD: TS. Lương Thái Ngọc
% Build:  .\build.ps1 slide-bao-ve.tex (chạy từ thư mục gốc)

\documentclass[fleqn]{beamer}
\usepackage[utf8]{inputenc}
\usepackage[T5]{fontenc}
\usepackage[vietnamese]{babel}
\usepackage{mathtools,amssymb,amsmath}
\usepackage{graphicx}
\usepackage{booktabs}
\usepackage{tikz}
\usetikzlibrary{shapes,arrows.meta,positioning}
\usepackage[backend=biber,style=numeric]{biblatex}
\addbibresource{../bao-cao/tai-lieu-tham-khao.bib}

\hypersetup{unicode}

\usetheme{Madrid}
\usecolortheme{default}
\graphicspath{{images/}}

\AtBeginSection[]{\begin{frame}\frametitle{Nội dung}\tableofcontents[currentsection]\end{frame}}

\title[Thiết kế Kiến trúc Mạng Không Dây DTHU]{THIẾT KẾ KIẾN TRÚC TỔNG THỂ HẠ TẦNG MẠNG KHÔNG DÂY TẠI TRƯỜNG ĐẠI HỌC ĐỒNG THÁP}
\author[Nguyễn Thị Kỳ Duyên]{Sinh viên: Nguyễn Thị Kỳ Duyên \\ \small MSSV: 0023410647 -- Lớp: ĐHCNTT23B-CS}
\institute[DTHU]{Khoa Công nghệ và Kỹ thuật, Trường Đại học Đồng Tháp\\ Ngành: Khoa học máy tính -- Chuyên ngành: Mạng Máy Tính và An Ninh\\ \textbf{GVHD: TS. Lương Thái Ngọc}}
\date{Tháng 04/2027}

\begin{document}

\begin{frame}\titlepage\end{frame}

\begin{frame}{Nội dung trình bày}
  \tableofcontents
\end{frame}

\section{Đặt vấn đề \& Mục tiêu đề tài}

\begin{frame}{Thực trạng hạ tầng Wi-Fi tại Trường Đại học Đồng Tháp}
  \begin{itemize}
    \item \textbf{Bối cảnh chuyển đổi số giáo dục đại học:}
      \begin{itemize}
        \item Quy mô hơn 15.000 sinh viên, cán bộ và giảng viên.
        \item Mỗi cá nhân sử dụng từ 1 đến 2 thiết bị di động (BYOD: laptop, smartphone, tablet).
        \item Nhu cầu truy cập tài liệu số, hệ thống E-learning, thi trực tuyến tăng vọt.
      \end{itemize}
    \item \textbf{Hạn chế của hạ tầng hiện hữu:}
      \begin{itemize}
        \item Phát triển tự phát theo nhu cầu cục bộ, thiếu quy hoạch tổng thể đồng bộ.
        \item Giờ cao điểm xảy ra hiện tượng sóng báo ``đầy vạch'' nhưng không thể truy cập Internet.
        \item Can nhiễu phổ vô tuyến nghiêm trọng giữa các Access Point chuẩn cũ (802.11n/ac).
      \end{itemize}
  \end{itemize}
\end{frame}

\begin{frame}{Mục tiêu và phạm vi nghiên cứu}
  \begin{columns}
    \begin{column}{0.52\textwidth}
      \textbf{Mục tiêu nghiên cứu:}
      \begin{itemize}\small
        \item Khảo sát, đo kiểm định lượng hiện trạng sóng Wi-Fi toàn campus DTHU.
        \item Chẩn đoán chính xác 3 điểm nghẽn cổ chai của hệ thống.
        \item Thiết kế kiến trúc mạng không dây chuẩn hóa (Wi-Fi 6 + WLC tập trung).
        \item Mô phỏng kiểm chứng trên \textbf{Cisco Packet Tracer} và đề xuất lộ trình tối ưu.
      \end{itemize}
    \end{column}
    \begin{column}{0.48\textwidth}
      \textbf{Phạm vi đề tài:}
      \begin{itemize}\small
        \item \textbf{Khu vực khảo sát:} Giảng đường A1, Hành chính B3, Căn tin B5, Thư viện, KTX B6, các tòa H1, H2, H3, C1, C2...
        \item \textbf{Công cụ đo kiểm:} Wi-Fi Analyzer, Speedtest, Wireshark.
        \item \textbf{Mô phỏng:} Cisco Packet Tracer 8.x.
        \item \textbf{Giới hạn:} Đánh giá giải pháp trên mô hình mô phỏng kỹ thuật, không can thiệp phần cứng thực tế.
      \end{itemize}
    \end{column}
  \end{columns}
\end{frame}

\section{Cơ sở lý thuyết mạng không dây}

\begin{frame}{Hệ thống chuẩn IEEE 802.11 và Công nghệ Wi-Fi 6}
  \begin{itemize}\small
    \item \textbf{Sự tiến hóa chuẩn IEEE 802.11:}
      \begin{itemize}\footnotesize
        \item 802.11b/g (2.4 GHz, 11/54 Mbps) $\rightarrow$ 802.11n (Wi-Fi 4, 600 Mbps) $\rightarrow$ 802.11ac (Wi-Fi 5, 5 GHz, 1.73 Gbps) $\rightarrow$ \textbf{802.11ax (Wi-Fi 6, 9.6 Gbps)}.
      \end{itemize}
    \item \textbf{Ba đột phá cốt lõi của Wi-Fi 6 cho môi trường mật độ cao (High Density):}
      \begin{itemize}\footnotesize
        \item \textbf{OFDMA (Orthogonal Frequency Division Multiple Access):} Chia nhỏ kênh thành các Resource Units (RU), phục vụ đồng thời nhiều thiết bị, triệt tiêu độ trễ hàng đợi.
        \item \textbf{MU-MIMO 8x8 (Uplink/Downlink):} Tăng gấp đôi luồng dữ liệu đồng thời so với Wi-Fi 5.
        \item \textbf{BSS Coloring (Tô màu khung tin):} Phân biệt sóng lân cận cùng tần số, giảm thiểu va chạm kênh truyền (Clear Channel Assessment - CCA).
      \end{itemize}
  \end{itemize}
\end{frame}

\begin{frame}{Mô hình quản lý tập trung WLC và Fast Roaming}
  \begin{columns}
    \begin{column}{0.5\textwidth}
      \textbf{Kiến trúc WLC (Split-MAC):}
      \begin{itemize}\small
        \item Phân tách chức năng giữa AP hạng nhẹ (Lightweight AP) và WLC thông qua giao thức \textbf{CAPWAP}.
        \item WLC đảm nhận: Quản lý sóng tự động (RRM), xác thực tập trung, cấu hình đồng bộ, QoS.
        \item Khắc phục hoàn toàn tình trạng AP độc lập cấu hình rời rạc.
      \end{itemize}
    \end{column}
    \begin{column}{0.5\textwidth}
      \textbf{Chuyển vùng nhanh (Fast Roaming):}
      \begin{itemize}\small
        \item \textbf{IEEE 802.11r (Fast Transition):} Giảm thời gian bắt tay mã hóa chuyển vùng xuống $< 50\text{ms}$.
        \item \textbf{IEEE 802.11k/v:} Hỗ trợ báo cáo lân cận và điều hướng sóng (BSS Transition), triệt tiêu lỗi thiết bị bám dính AP ở xa (\textit{Sticky Client}).
      \end{itemize}
    \end{column}
  \end{columns}
\end{frame}

\section{Khảo sát \& Đánh giá hiện trạng DTHU}

\begin{frame}{Khảo sát thực địa hạ tầng thiết bị hiện hữu}
  \begin{itemize}
    \item \textbf{Bộ điều khiển trung tâm (Wireless Switch):} Motorola RFS6000.
      \begin{itemize}\small
        \item Năng lực tối đa: Chỉ cấp phép cho tối đa \textbf{2.000 kết nối đồng thời}.
        \item Giờ cao điểm (09h30--10h30, 14h00--15h30): Lưu lượng thực tế đạt \textbf{2.150 -- 2.300 thiết bị} $\rightarrow$ Tràn bảng cấp phát DHCP, từ chối kết nối mới.
      \end{itemize}
    \item \textbf{Mạng lưới Access Point tại các khối nhà:}
      \begin{itemize}\small
        \item Các khối giảng đường lớn (Tòa A1, Tòa C2): Mật độ thiết bị dày đặc nhưng lắp đặt AP chuẩn cũ.
        \item Tòa hành chính (B1, B2, B3): Tường bê tông dày, suy hao tín hiệu cao.
        \item Không gian mở (Căn tin B5): Tín hiệu lan tỏa rộng nhưng giao thoa nhiễu mạnh.
      \end{itemize}
  \end{itemize}
\end{frame}

\begin{frame}{Số liệu đo kiểm thực tế (Wi-Fi Analyzer \& Speedtest)}
  \begin{table}\small\centering
  \begin{tabular}{lcccc}
    \toprule
    \textbf{Khu vực} & \textbf{Khung giờ} & \textbf{Download} & \textbf{Upload} & \textbf{Độ trễ Ping} \\
    \midrule
    Tòa A1 (Giảng đường) & 6h00 (vắng) & 20.08 Mbps & 18.02 Mbps & 23.45 ms \\
                         & 9h00 (cao điểm) & \textbf{5.19 Mbps} & \textbf{6.98 Mbps} & \textbf{110.0 ms} \\
    \midrule
    Tòa B3 (Hành chính)  & 7h00 (vắng) & 23.89 Mbps & 20.18 Mbps & 28.90 ms \\
                         & 10h00 (cao điểm)& 15.35 Mbps & 16.45 Mbps & 130.0 ms \\
    \midrule
    Tòa B5 (Căn tin mở)  & 7h20 (vắng) & 19.80 Mbps & 18.89 Mbps & 20.05 ms \\
                         & 11h30 (cao điểm)& 9.09 Mbps & 10.29 Mbps & 100.0 ms \\
    \bottomrule
  \end{tabular}
  \end{table}
  \vspace{0.15cm}
  $\Rightarrow$ Tốc độ tải xuống tại Giảng đường A1 giờ cao điểm sụt giảm \textbf{gần 74\%}, Ping vượt ngưỡng 100ms.
\end{frame}

\begin{frame}{Ba điểm nghẽn kỹ thuật cốt lõi}
  \begin{enumerate}
    \item \textbf{Quá tải phần cứng WLC:} Thiết bị RFS6000 chạm trần 2.000 thiết bị, CPU nghẽn và kiệt tài nguyên cấp phát DHCP.
    \item \textbf{Can nhiễu đồng kênh (Co-channel Interference):}
      \begin{itemize}\small
        \item Băng tần 2.4 GHz bị lạm dụng, cấu hình kênh chồng lấn thay vì cố định (1, 6, 11).
        \item Tranh chấp thời gian phát sóng (\textit{Airtime Contention}) làm giảm nghiêm trọng hiệu năng truyền tải.
      \end{itemize}
    \item \textbf{Thiết bị phát sóng ngoại lai (Rogue AP):}
      \begin{itemize}\small
        \item Xuất hiện nhiều router Wi-Fi cá nhân cắm sai quy định (ví dụ SSID: TP-LINK\_E9A6).
        \item Cấp phát DHCP giả mạo làm mất mạng diện rộng và tiềm ẩn nguy cơ bảo mật.
      \end{itemize}
  \end{enumerate}
\end{frame}

\section{Thiết kế kiến trúc tổng thể \& Giải pháp}

\begin{frame}{Mô hình thiết kế kiến trúc phân cấp 3 lớp}
  \begin{itemize}\small
    \item \textbf{Lớp Lõi (Core Layer):} Cặp Switch Core hiệu năng cao hoạt động chế độ dự phòng chủ động, băng thông chuyển mạch tốc độ cao 10/40 Gbps.
    \item \textbf{Lớp Phân phối (Distribution Layer):} Định tuyến đa WAN, cân bằng tải động cho khối Sinh viên và chính sách định tuyến tĩnh PBR ưu tiên cho khối Giảng viên/Hành chính.
    \item \textbf{Lớp Truy cập (Access Layer):} Hệ thống Switch Access chuẩn \textbf{PoE+ (802.3at)} cấp nguồn và truyền dẫn GigE tới các AP.
    \item \textbf{Hạ tầng Điều khiển trung tâm:} Nâng cấp WLC năng lực quản trị \textbf{5.000 -- 10.000 kết nối đồng thời}, hỗ trợ cơ chế dự phòng $N+1$ High Availability.
  \end{itemize}
\end{frame}

\begin{frame}{Quy hoạch VLAN, Dải IP và Chiến lược DHCP}
  \begin{table}\small\centering
  \begin{tabular}{cllll}
    \toprule
    \textbf{VLAN} & \textbf{Đối tượng} & \textbf{Dải mạng IP (Subnet)} & \textbf{SSID} & \textbf{Lease Time} \\
    \midrule
    10 & Cán bộ / GV & 192.168.10.0/24 (254 IP) & DTHU-STAFF & 24 giờ \\
    20 & Sinh viên Tòa A & 10.10.0.0/20 (4.094 IP) & DTHU-STUDENT & 2 -- 4 giờ \\
    30 & Sinh viên Tòa B,C & 10.20.0.0/20 (4.094 IP) & DTHU-STUDENT & 2 -- 4 giờ \\
    40 & Khách vãng lai & 172.16.40.0/24 (254 IP) & DTHU-GUEST & 1 giờ \\
    \bottomrule
  \end{tabular}
  \end{table}
  \vspace{0.2cm}
  \begin{itemize}\small
    \item \textbf{Đột phá giải quyết cạn kiệt IP:} Áp dụng subnet /20 cung cấp hơn 8.000 IP cho sinh viên kết hợp giảm Lease Time xuống 2--4h để thu hồi IP nhanh chóng.
    \item \textbf{Bảo mật:} Cách ly máy trạm lớp 2 (Layer 2 Client Isolation), ngăn chặn quét mạng nội bộ và lây lan mã độc.
  \end{itemize}
\end{frame}

\begin{frame}{Giải pháp Tối ưu Vô tuyến và An ninh mạng}
  \begin{itemize}
    \item \textbf{Quy hoạch AP chuẩn Wi-Fi 6 tại điểm nóng:}
      \begin{itemize}\small
        \item Trang bị AP Wi-Fi 6 (OFDMA, MU-MIMO) tại Giảng đường A1, Thư viện, Căn tin B5.
        \item Tái sử dụng các AP Wi-Fi 4/5 còn tốt cho các phòng làm việc nhỏ ít người.
      \end{itemize}
    \item \textbf{Quản lý tài nguyên vô tuyến thông minh (RRM):}
      \begin{itemize}\small
        \item Tự động gán kênh sạch (DCA) và tự cân chỉnh công suất phát sóng (TPC).
      \end{itemize}
    \item \textbf{Xác thực tập trung cấp Doanh nghiệp (Enterprise Security):}
      \begin{itemize}\small
        \item Chuyển đổi từ mật khẩu chia sẻ (WPA2-PSK) sang xác thực định danh cá nhân \textbf{WPA2/WPA3-Enterprise (802.1X/EAP)} tích hợp máy chủ RADIUS/LDAP.
        \item Hệ thống tự động dò quét và vô hiệu hóa thiết bị ngoại lai Rogue AP (WIDS/WIPS).
      \end{itemize}
  \end{itemize}
\end{frame}

\section{Kết quả mô phỏng \& Đánh giá}

\begin{frame}{Mô hình mô phỏng trên Cisco Packet Tracer}
  \begin{itemize}
    \item \textbf{Xây dựng kiến trúc hoàn chỉnh:}
      \begin{itemize}\small
        \item Tích hợp Router Biên $\rightarrow$ Core Switch $\rightarrow$ Distribution Switches $\rightarrow$ PoE Access Switches.
        \item WLC kết nối máy chủ dịch vụ RADIUS Server, DHCP Server, DNS và Web Server.
        \item Cấu hình AP phát đồng thời đa SSID với các chính sách VLAN độc lập.
      \end{itemize}
    \item \textbf{Kiểm thử các kịch bản thực tế:}
      \begin{itemize}\small
        \item Kịch bản cấp phát đồng thời lượng lớn IP qua relay agent.
        \item Kịch bản cô lập Client-to-Client trong mạng sinh viên.
        \item Kịch bản chuyển vùng mượt mà giữa các AP khi di chuyển.
      \end{itemize}
  \end{itemize}
\end{frame}

\begin{frame}{So sánh định lượng trước và sau tối ưu}
  \begin{table}\small\centering
  \begin{tabular}{lll}
    \toprule
    \textbf{Chỉ số kỹ thuật} & \textbf{Hiện trạng khảo sát} & \textbf{Kỳ vọng sau tối ưu} \\
    \midrule
    Năng lực chịu tải đồng thời & $\sim 2.000$ thiết bị (Gây nghẽn) & \textbf{$> 10.000$ thiết bị} \\
    Nhiễu đồng kênh Tòa A1 & Rất cao (Lỗi quy hoạch) & \textbf{Thấp} (Thuật toán DCA/RRM) \\
    Độ trễ Ping giờ cao điểm & $> 100$ ms & \textbf{$< 30$ ms} (Wi-Fi 6 OFDMA) \\
    Bảo mật xác thực & PSK dùng chung (Rủi ro cao) & \textbf{802.1X} (Định danh cá nhân) \\
    Thời gian chuyển vùng & Kém (Thường bị ngắt kết nối) & \textbf{$< 50$ ms} (IEEE 802.11r/k/v) \\
    Tỷ lệ mất gói tin & 7\% -- 12\% & \textbf{$< 2\%$} \\
    \bottomrule
  \end{tabular}
  \end{table}
\end{frame}

\section{Kết luận \& Hướng phát triển}

\begin{frame}{Kết luận và Đóng góp của đề tài}
  \begin{itemize}
    \item \textbf{Kết quả đạt được:}
      \begin{itemize}\small
        \item Hoàn thành 100\% các mục tiêu trong Đề cương nghiên cứu đã được phê duyệt.
        \item Cung cấp bộ dữ liệu khảo sát thực địa toàn diện và bóc tách chính xác nguyên nhân nghẽn Wi-Fi tại DTHU.
        \item Xây dựng thành công bản vẽ kiến trúc mạng không dây tổng thể đáp ứng tăng trưởng 5 năm tới.
      \end{itemize}
    \item \textbf{Ý nghĩa thực tiễn:}
      \begin{itemize}\small
        \item Làm tài liệu kỹ thuật tham khảo giá trị cho Trung tâm CNTT - Trường Đại học Đồng Tháp.
        \item Phương án có tính khả thi cao, tiết kiệm ngân sách nhờ tái sử dụng thiết bị cũ và tận dụng máy chủ LDAP sẵn có.
      \end{itemize}
    \item \textbf{Hướng phát triển tương lai:}
      \begin{itemize}\small
        \item Triển khai thí điểm (PoC) tại Giảng đường A1 để đo kiểm thực nghiệm.
        \item Tích hợp AI/ML trong vận hành mạng tự động (AIOps).
        \item Nghiên cứu chuẩn bị hạ tầng đón đầu tiêu chuẩn Wi-Fi 7 (802.11be).
      \end{itemize}
  \end{itemize}
\end{frame}

\begin{frame}{Tài liệu tham khảo chính}
  \nocite{*}
  \printbibliography
\end{frame}

\begin{frame}
  \centering
  \Large\bfseries
  Xin chân thành cảm ơn quý Thầy, Cô\\ trong Hội đồng đã chú ý lắng nghe!\\[1cm]
  \normalsize Kính mong nhận được ý kiến đóng góp quý báu của quý Thầy, Cô.
\end{frame}

\end{document}
